% Complete reviewed Unit 5; editorial master: chapters/ch05/answer-key.md
% Markdown SHA-256: 16d3c595ef46d19d7f79a32b94063247effa7793f8f42feec4ad5b69ce4a5c70
% One-use initial assembly followed by complete content and page review.
\newcommand{\runningtitle}{Unit 5. Answer key}
% Layout inherited from the reviewed previous course; set a running title per unit.
% This is a small hand-maintained TeX layout, not a conversion pipeline.
\RequirePackage{fix-cm}
\documentclass[11pt,twoside]{article}
\usepackage{fontspec}
\setmainfont{Cambria}[
  Renderer=HarfBuzz,
  BoldFont=cambriab.ttf,
  ItalicFont=cambriai.ttf,
  BoldItalicFont=cambriaz.ttf
]
\setmonofont{lmmono10-regular.otf}[BoldFont=lmmonolt10-bold.otf]
\usepackage[paperwidth=165mm,paperheight=241mm,inner=21mm,outer=19mm,
  top=22mm,bottom=20mm,headheight=14pt,headsep=7mm,footskip=10mm]{geometry}
\usepackage{xcolor,array,booktabs,longtable,tabularx}
\usepackage{enumitem,fancyhdr}
\usepackage[unicode,hidelinks]{hyperref}
\definecolor{ink}{HTML}{000000}
\definecolor{muted}{HTML}{4C5558}
\setlength{\parindent}{0pt}
\setlength{\parskip}{3.5pt}
\AtBeginDocument{\fontsize{11}{14.3}\selectfont}
\setlength{\tabcolsep}{4pt}
\renewcommand{\arraystretch}{1.10}
\setlist[itemize]{leftmargin=1.3em,label=\textbullet,itemsep=2pt,topsep=3pt,parsep=0pt}
\setlist[enumerate]{leftmargin=1.8em,itemsep=3pt,topsep=4pt,parsep=0pt}
\widowpenalty=10000
\clubpenalty=10000
\emergencystretch=1em
\raggedbottom
\pagestyle{fancy}
\fancyhf{}
\providecommand{\runningtitle}{Old English}
\fancyhead[LE,RO]{\small\itshape \runningtitle}
\fancyhead[LO,RE]{\small Old English}
\fancyfoot[LE,RO]{\small\thepage}
\renewcommand{\headrulewidth}{0.3pt}
\newcommand{\booktitle}[2]{
  {\fontsize{9}{11}\selectfont\color{muted}#1\par}
  \vspace{3pt}
  {\fontsize{20}{24}\selectfont\color{ink}#2\par}
  \vspace{6pt}
}
\newcommand{\source}[1]{\par{\footnotesize\color{muted}#1\par}}
\newcommand{\readingtext}[1]{
  \par\addvspace{5pt}
  \begingroup\leftskip=1.5em
    \fontsize{12}{16}\selectfont #1\par
  \endgroup\addvspace{5pt}
}
\makeatletter
\renewcommand{\section}{\@startsection{section}{1}{\z@}
  {-2.4ex \@plus -0.6ex \@minus -0.2ex}{0.8ex \@plus 0.2ex}
  {\normalfont\fontsize{12}{15}\selectfont\bfseries\color{ink}}}
\renewcommand{\subsection}{\@startsection{subsection}{2}{\z@}
  {-1.8ex \@plus -0.5ex \@minus -0.2ex}{0.5ex \@plus 0.2ex}
  {\normalfont\fontsize{11}{14}\selectfont\bfseries\color{ink}}}
\makeatother

\hyphenpenalty=10000
\exhyphenpenalty=10000
\emergencystretch=3em
\newcommand{\itempara}[2]{\par\addvspace{3pt}\noindent\hangindent=1.8em\hangafter=1\makebox[1.8em][l]{#1.}#2\par}
\newcommand{\reserve}[1]{\par\penalty-100\begingroup\dimen0=\pagegoal\advance\dimen0 by-\pagetotal\ifdim\dimen0<#1\newpage\fi\endgroup}
\begin{document}
\addfontfeatures{Ligatures=TeXOff}
\raggedright
\booktitle{UNIT 5}{Answer key}

\reserve{7\baselineskip}
% mdparagraph:P001:begin
\section*{Using the key}
% mdparagraph:P001:end

% mdparagraph:P002:begin
The models use the lesson's selected forms. In a form exercise, the target is the stated verb, tense, person and number; a documented variant can be valid Old English while differing from the selected table. In a translation, preserve the people, number, tense, objects, possessors and other stated relations. Word order may vary when those relations remain clear. The notes identify the main targets and likely errors; they do not make each model the only acceptable string.
% mdparagraph:P002:end

% mdparagraph:P003:begin
The present second and third singular often change the stem vowel or consonants. The preterite second singular uses the plural-preterite stem plus \textbf{-e}; all three plural persons use \textbf{-on}. A dual subject also takes a plural verb. Preserve quantity: \textbf{æ/\allowbreak{}ǣ}, \textbf{ea/\allowbreak{}ēa}, \textbf{ie/\allowbreak{}īe} and \textbf{e/\allowbreak{}ē} are not interchangeable. All requested finite forms below are indicative. [1]
% mdparagraph:P003:end

% mdparagraph:P004:begin
Infinitives and participles printed in four-part entries remain supplied citations. Identifying a citation does not require producing an infinitive clause, perfect or passive. Likewise, the reading's local help remains available for wishes, naming, relative and content clauses, passive or resultative phrases, and the other supplied uses. Exercise I.5.6 supplies its whole \textbf{eall ðā ðing} phrase; no independent quantifier production is required.
% mdparagraph:P004:end

% mdparagraph:P005:begin
Both \textbf{sēon/\allowbreak{}gesēon} are separately taught for ordinary seeing, and both \textbf{scieppan/\allowbreak{}gescieppan} for creating. Either member can satisfy an unqualified English prompt in the appropriate finite form. A task that explicitly names one verb still has that lexical target. These checked pairs do not authorize arbitrary prefix addition or removal. The answers and translations are newly prepared teaching models, not a transcription of Cook's answer key.
% mdparagraph:P005:end

\reserve{7\baselineskip}
% mdparagraph:P006:begin
\section*{Principal parts}
% mdparagraph:P006:end

\reserve{9\baselineskip}
\par\begin{minipage}{\linewidth}\raggedright
% mdparagraph:P007:begin
\textbf{P1.}
% mdparagraph:P007:end

% item: NEW-U05-P01
% key: KEY-NEW-U05-P01
% mdtable:T01:begin
\par\begingroup
\fontsize{10.5}{13.2}\selectfont
\setlength{\LTleft}{0pt}\setlength{\LTright}{\fill}
\noindent
\begin{tabular}{@{}>{\raggedright\arraybackslash}p{22.17mm}>{\raggedright\arraybackslash}p{22.17mm}>{\raggedright\arraybackslash}p{74.04mm}@{}}
\toprule
\textbf{Position} & \textbf{Form} & \textbf{Function in the entry} \\
\midrule
First & beran & Infinitive citation \\
Second & bær & Preterite singular principal part: first or third person \\
Third & bǣron & Preterite plural principal part: first, second or third person \\
Fourth & boren & Past participle citation \\
\bottomrule
\end{tabular}
\par\endgroup

% mdtable:T01:end
% mdparagraph:P008:begin
\textbf{Bær} and \textbf{bǣron} are finite indicative forms. \textbf{Beran} and \textbf{boren} are nonfinite citations and do not specify indicative person and number. The entry distinguishes the singular and plural past stems; neither the infinitive ending \textbf{-an} nor participial \textbf{-en} is another personal ending. [1]
% mdparagraph:P008:end

\end{minipage}\par\addvspace{4pt}

\par\begin{minipage}{\linewidth}\raggedright
% mdparagraph:P009:begin
\textbf{P2.}
% mdparagraph:P009:end

% item: NEW-U05-P02
% key: KEY-NEW-U05-P02
% mdparagraph:P010:begin
\textbf{Geaf}, from \textbf{giefan}, is strong Class V: the preterite is marked by the changed stem vowel, without a dental preterite suffix. \textbf{Sealde}, from \textbf{sęllan}, is weak: it contains the dental preterite marker \textbf{-d-}. Each can be first- or third-person singular preterite indicative when given without a subject. Their common English meaning, “gave,” does not determine their conjugation.
% mdparagraph:P010:end

\end{minipage}\par\addvspace{4pt}

\par\begin{minipage}{\linewidth}\raggedright
% mdparagraph:P011:begin
\textbf{P3.}
% mdparagraph:P011:end

% item: NEW-U05-P03
% key: KEY-NEW-U05-P03
% mdtable:T02:begin
\par\begingroup
\fontsize{10.5}{13.2}\selectfont
\setlength{\LTleft}{0pt}\setlength{\LTright}{\fill}
\noindent
\begin{tabular}{@{}>{\raggedright\arraybackslash}p{57.78mm}>{\raggedright\arraybackslash}p{63.41mm}@{}}
\toprule
\textbf{Subject} & \textbf{Preterite form} \\
\midrule
I & bær \\
You, singular & bǣre \\
He & bær \\
We & bǣron \\
\bottomrule
\end{tabular}
\par\endgroup

% mdtable:T02:end
% mdparagraph:P012:begin
\textbf{Bǣre} and \textbf{bǣron} use the plural-preterite stem \textbf{bǣr-}. First and third singular use short \textbf{æ} in \textbf{bær}. A second-singular answer made from that short-vowel stem misses both the required stem and the distinction in quantity.
% mdparagraph:P012:end

\end{minipage}\par\addvspace{4pt}

\par\begin{minipage}{\linewidth}\raggedright
% mdparagraph:P013:begin
\textbf{P4.}
% mdparagraph:P013:end

% item: NEW-U05-P04
% key: KEY-NEW-U05-P04
% mdparagraph:P014:begin
“You, singular, created” is \textbf{gescēope}; “we created” is \textbf{gescēopon}. Both use the supplied plural stem \textbf{gescēop-}, with \textbf{-e} and \textbf{-on} respectively. \textbf{Gesceapen} is the supplied past participle, not a finite preterite. Its final \textbf{-en} does not identify a plural subject.
% mdparagraph:P014:end

% mdparagraph:P015:begin
The separately documented \textbf{ō} series is valid, but this task asks the learner to use the printed \textbf{ēo} entry. Its nonfinite citations may remain in view. [4]
% mdparagraph:P015:end

\end{minipage}\par\addvspace{4pt}

\reserve{7\baselineskip}
% mdparagraph:P016:begin
\section*{Finite forms}
% mdparagraph:P016:end

% mdparagraph:P017:begin
For F1–F24, each table gives exactly the requested cells. The explanation identifies the relevant stem or ending; the shared notes above apply to all rows. A response in another person, number or tense changes the requested grammatical meaning even when it is a valid form elsewhere in the paradigm.
% mdparagraph:P017:end

\reserve{6\baselineskip}
\par\begin{minipage}{\linewidth}\raggedright
% mdparagraph:P018:begin
\textbf{F1.}
% mdparagraph:P018:end

% item: NEW-U05-F01
% key: KEY-NEW-U05-F01
% mdtable:T03:begin
\par\begingroup
\fontsize{10.5}{13.2}\selectfont
\setlength{\LTleft}{0pt}\setlength{\LTright}{\fill}
\noindent
\begin{tabular}{@{}>{\raggedright\arraybackslash}p{37.46mm}>{\raggedright\arraybackslash}p{40.46mm}>{\raggedright\arraybackslash}p{40.46mm}@{}}
\toprule
\textbf{Present 2sg} & \textbf{Preterite 1sg} & \textbf{Preterite 2sg} \\
\midrule
birst & bær & bǣre \\
\bottomrule
\end{tabular}
\par\endgroup

% mdtable:T03:end
% mdparagraph:P019:begin
The present has selected \textbf{i}; \textbf{bierst} is a documented alternative. The past second singular has long \textbf{ǣ}, unlike first-singular \textbf{bær}. [2]
% mdparagraph:P019:end

\end{minipage}\par\addvspace{4pt}

\par\begin{minipage}{\linewidth}\raggedright
% mdparagraph:P020:begin
\textbf{F2.}
% mdparagraph:P020:end

% item: NEW-U05-F02
% key: KEY-NEW-U05-F02
% mdtable:T04:begin
\par\begingroup
\fontsize{10.5}{13.2}\selectfont
\setlength{\LTleft}{0pt}\setlength{\LTright}{\fill}
\noindent
\begin{tabular}{@{}>{\raggedright\arraybackslash}p{37.46mm}>{\raggedright\arraybackslash}p{40.46mm}>{\raggedright\arraybackslash}p{40.46mm}@{}}
\toprule
\textbf{Present 3sg} & \textbf{Preterite 3sg} & \textbf{Preterite 1pl} \\
\midrule
nimð & nōm & nōmon \\
\bottomrule
\end{tabular}
\par\endgroup

% mdtable:T04:end
% mdparagraph:P021:begin
The selected past has long \textbf{ō} in both singular and plural. It does not use the ordinary \textbf{bær/\allowbreak{}bǣron} vowel series. Documented \textbf{nam/\allowbreak{}nāmon} belong to an alternative past series and differ from the selected target. [2]
% mdparagraph:P021:end

\end{minipage}\par\addvspace{4pt}

\par\begin{minipage}{\linewidth}\raggedright
% mdparagraph:P022:begin
\textbf{F3.}
% mdparagraph:P022:end

% item: NEW-U05-F03
% key: KEY-NEW-U05-F03
% mdtable:T05:begin
\par\begingroup
\fontsize{10.5}{13.2}\selectfont
\setlength{\LTleft}{0pt}\setlength{\LTright}{\fill}
\noindent
\begin{tabular}{@{}>{\raggedright\arraybackslash}p{28.32mm}>{\raggedright\arraybackslash}p{28.32mm}>{\raggedright\arraybackslash}p{28.32mm}>{\raggedright\arraybackslash}p{30.59mm}@{}}
\toprule
\textbf{Present 1sg} & \textbf{Present 2sg} & \textbf{Present 3sg} & \textbf{Preterite 3pl} \\
\midrule
cume & cymst & cymð & cōmon \\
\bottomrule
\end{tabular}
\par\endgroup

% mdtable:T05:end
% mdparagraph:P023:begin
Present \textbf{u} becomes \textbf{y} in the second and third singular. The preterite has long \textbf{ō}. \textbf{Cwōmon} is a documented alternative to selected \textbf{cōmon}; it is not a reason to insert \textbf{w} in every form. [2]
% mdparagraph:P023:end

\end{minipage}\par\addvspace{4pt}

\par\begin{minipage}{\linewidth}\raggedright
% mdparagraph:P024:begin
\textbf{F4.}
% mdparagraph:P024:end

% item: NEW-U05-F04
% key: KEY-NEW-U05-F04
% mdtable:T06:begin
\par\begingroup
\fontsize{10.5}{13.2}\selectfont
\setlength{\LTleft}{0pt}\setlength{\LTright}{\fill}
\noindent
\begin{tabular}{@{}>{\raggedright\arraybackslash}p{38.43mm}>{\raggedright\arraybackslash}p{38.43mm}>{\raggedright\arraybackslash}p{41.51mm}@{}}
\toprule
\textbf{Present 2sg} & \textbf{Present 3sg} & \textbf{Preterite 2sg} \\
\midrule
spricst & spricð & sprǣce \\
\bottomrule
\end{tabular}
\par\endgroup

% mdtable:T06:end
% mdparagraph:P025:begin
The present second and third singular have \textbf{i}. The past second singular uses the long-\textbf{ǣ} plural stem, not the short-\textbf{æ} singular principal part \textbf{spræc}. [3]
% mdparagraph:P025:end

\end{minipage}\par\addvspace{4pt}

\par\begin{minipage}{\linewidth}\raggedright
% mdparagraph:P026:begin
\textbf{F5.}
% mdparagraph:P026:end

% item: NEW-U05-F05
% key: KEY-NEW-U05-F05
% mdtable:T07:begin
\par\begingroup
\fontsize{10.5}{13.2}\selectfont
\setlength{\LTleft}{0pt}\setlength{\LTright}{\fill}
\noindent
\begin{tabular}{@{}>{\raggedright\arraybackslash}p{38.43mm}>{\raggedright\arraybackslash}p{38.43mm}>{\raggedright\arraybackslash}p{41.51mm}@{}}
\toprule
\textbf{Present 2sg} & \textbf{Present 3sg} & \textbf{Preterite 3pl} \\
\midrule
cwist & cwið & cwǣdon \\
\bottomrule
\end{tabular}
\par\endgroup

% mdtable:T07:end
% mdparagraph:P027:begin
The second-singular present loses stem \textbf{ð} before \textbf{-st}. The plural preterite has \textbf{d}, together with long \textbf{ǣ}. Do not carry the first/\allowbreak{}third-singular \textbf{cwæð} stem unchanged into the plural. [3]
% mdparagraph:P027:end

\end{minipage}\par\addvspace{4pt}

\par\begin{minipage}{\linewidth}\raggedright
% mdparagraph:P028:begin
\textbf{F6.}
% mdparagraph:P028:end

% item: NEW-U05-F06
% key: KEY-NEW-U05-F06
% mdtable:T08:begin
\par\begingroup
\fontsize{10.5}{13.2}\selectfont
\setlength{\LTleft}{0pt}\setlength{\LTright}{\fill}
\noindent
\begin{tabular}{@{}>{\raggedright\arraybackslash}p{27.78mm}>{\raggedright\arraybackslash}p{27.78mm}>{\raggedright\arraybackslash}p{30.0mm}>{\raggedright\arraybackslash}p{30.0mm}@{}}
\toprule
\textbf{Present 2sg} & \textbf{Present 3sg} & \textbf{Preterite 1sg} & \textbf{Preterite 1pl} \\
\midrule
giefst & giefð & geaf & gēafon \\
\bottomrule
\end{tabular}
\par\endgroup

% mdtable:T08:end
% mdparagraph:P029:begin
The selected present retains \textbf{ie}. Short \textbf{ea} in singular \textbf{geaf} contrasts with long \textbf{ēa} in plural \textbf{gēafon}; the \textbf{g}-initial spelling does not remove that quantity distinction. [3]
% mdparagraph:P029:end

\end{minipage}\par\addvspace{4pt}

\par\begin{minipage}{\linewidth}\raggedright
% mdparagraph:P030:begin
\textbf{F7.}
% mdparagraph:P030:end

% item: NEW-U05-F07
% key: KEY-NEW-U05-F07
% mdtable:T09:begin
\par\begingroup
\fontsize{10.5}{13.2}\selectfont
\setlength{\LTleft}{0pt}\setlength{\LTright}{\fill}
\noindent
\begin{tabular}{@{}>{\raggedright\arraybackslash}p{38.43mm}>{\raggedright\arraybackslash}p{38.43mm}>{\raggedright\arraybackslash}p{41.51mm}@{}}
\toprule
\textbf{Present 1sg} & \textbf{Present 3sg} & \textbf{Preterite 2sg} \\
\midrule
bidde & bitt & bǣde \\
\bottomrule
\end{tabular}
\par\endgroup

% mdtable:T09:end
% mdparagraph:P031:begin
The first singular preserves \textbf{dd}. The third singular has the selected dental outcome \textbf{tt}; documented \textbf{bit} is also valid. Past \textbf{bǣde} uses the long-\textbf{ǣ} stem and is strong, despite its stem-final \textbf{d}. [3]
% mdparagraph:P031:end

\end{minipage}\par\addvspace{4pt}

\par\begin{minipage}{\linewidth}\raggedright
% mdparagraph:P032:begin
\textbf{F8.}
% mdparagraph:P032:end

% item: NEW-U05-F08
% key: KEY-NEW-U05-F08
% mdtable:T10:begin
\par\begingroup
\fontsize{10.5}{13.2}\selectfont
\setlength{\LTleft}{0pt}\setlength{\LTright}{\fill}
\noindent
\begin{tabular}{@{}>{\raggedright\arraybackslash}p{38.43mm}>{\raggedright\arraybackslash}p{38.43mm}>{\raggedright\arraybackslash}p{41.51mm}@{}}
\toprule
\textbf{Present 2sg} & \textbf{Present 3sg} & \textbf{Preterite 1pl} \\
\midrule
ligst & ligð & lǣgon \\
\bottomrule
\end{tabular}
\par\endgroup

% mdtable:T10:end
% mdparagraph:P033:begin
The selected present has \textbf{g}, not the \textbf{cg} of the citation \textbf{licgan}. Cook's \textbf{līst/\allowbreak{}līð} are documented contracted alternatives, with long \textbf{ī}. Wright also records past plural \textbf{lāgon}; the selected table uses \textbf{lǣgon}. [3]
% mdparagraph:P033:end

\end{minipage}\par\addvspace{4pt}

\par\begin{minipage}{\linewidth}\raggedright
% mdparagraph:P034:begin
\textbf{F9.}
% mdparagraph:P034:end

% item: NEW-U05-F09
% key: KEY-NEW-U05-F09
% mdtable:T11:begin
\par\begingroup
\fontsize{10.5}{13.2}\selectfont
\setlength{\LTleft}{0pt}\setlength{\LTright}{\fill}
\noindent
\begin{tabular}{@{}>{\raggedright\arraybackslash}p{37.46mm}>{\raggedright\arraybackslash}p{40.46mm}>{\raggedright\arraybackslash}p{40.46mm}@{}}
\toprule
\textbf{Present 3sg} & \textbf{Preterite 1sg} & \textbf{Preterite 2sg} \\
\midrule
sitt & sæt & sǣte \\
\bottomrule
\end{tabular}
\par\endgroup

% mdtable:T11:end
% mdparagraph:P035:begin
Present \textbf{sit} is a documented alternative to selected \textbf{sitt}. The past first singular has short \textbf{æ}; second-singular \textbf{sǣte} uses the plural stem with long \textbf{ǣ}. [3]
% mdparagraph:P035:end

\end{minipage}\par\addvspace{4pt}

\par\begin{minipage}{\linewidth}\raggedright
% mdparagraph:P036:begin
\textbf{F10.}
% mdparagraph:P036:end

% item: NEW-U05-F10
% key: KEY-NEW-U05-F10
% mdtable:T12:begin
\par\begingroup
\fontsize{10.5}{13.2}\selectfont
\setlength{\LTleft}{0pt}\setlength{\LTright}{\fill}
\noindent
\begin{tabular}{@{}>{\raggedright\arraybackslash}p{27.78mm}>{\raggedright\arraybackslash}p{27.78mm}>{\raggedright\arraybackslash}p{30.0mm}>{\raggedright\arraybackslash}p{30.0mm}@{}}
\toprule
\textbf{Present 2sg} & \textbf{Present 3sg} & \textbf{Preterite 1sg} & \textbf{Preterite 2sg} \\
\midrule
itst & itt & ǣt & ǣte \\
\bottomrule
\end{tabular}
\par\endgroup

% mdtable:T12:end
% mdparagraph:P037:begin
The present has \textbf{i} and dental combinations. Selected past \textbf{ǣt} is long even in the singular; do not impose the short-vowel pattern of \textbf{sæt} on this verb. Short \textbf{æt} and third-singular present \textbf{it} are documented alternatives in their respective positions. [3]
% mdparagraph:P037:end

\end{minipage}\par\addvspace{4pt}

\par\begin{minipage}{\linewidth}\raggedright
% mdparagraph:P038:begin
\textbf{F11.}
% mdparagraph:P038:end

% item: NEW-U05-F11
% key: KEY-NEW-U05-F11
% mdtable:T13:begin
\par\begingroup
\fontsize{10.5}{13.2}\selectfont
\setlength{\LTleft}{0pt}\setlength{\LTright}{\fill}
\noindent
\begin{tabular}{@{}>{\raggedright\arraybackslash}p{39.46mm}>{\raggedright\arraybackslash}p{39.46mm}>{\raggedright\arraybackslash}p{39.46mm}@{}}
\toprule
\textbf{Present 1sg} & \textbf{Present 2sg} & \textbf{Present 3pl} \\
\midrule
sēo & siehst & sēoð \\
\bottomrule
\end{tabular}
\par\endgroup

% mdtable:T13:end
% mdparagraph:P039:begin
The contracted first singular and plural have \textbf{ēo}, without an added ordinary \textbf{-e/\allowbreak{}-að}. Second-singular \textbf{siehst} has short \textbf{ie}. Long \textbf{īe} would confuse this selected seeing system with a different vowel pattern. [3]
% mdparagraph:P039:end

\end{minipage}\par\addvspace{4pt}

\par\begin{minipage}{\linewidth}\raggedright
% mdparagraph:P040:begin
\textbf{F12.}
% mdparagraph:P040:end

% item: NEW-U05-F12
% key: KEY-NEW-U05-F12
% mdtable:T14:begin
\par\begingroup
\fontsize{10.5}{13.2}\selectfont
\setlength{\LTleft}{0pt}\setlength{\LTright}{\fill}
\noindent
\begin{tabular}{@{}>{\raggedright\arraybackslash}p{37.46mm}>{\raggedright\arraybackslash}p{40.46mm}>{\raggedright\arraybackslash}p{40.46mm}@{}}
\toprule
\textbf{Present 3sg} & \textbf{Preterite 2sg} & \textbf{Preterite 2pl} \\
\midrule
gesiehð & gesāwe & gesāwon \\
\bottomrule
\end{tabular}
\par\endgroup

% mdtable:T14:end
% mdparagraph:P041:begin
Present \textbf{ie} is short. The plural-based past has long \textbf{ā} and \textbf{w}; \textbf{geseah} is the first/\allowbreak{}third singular past and does not supply the requested second-singular stem. Keep the separately learned \textbf{ge-} of the named verb. [3]
% mdparagraph:P041:end

\end{minipage}\par\addvspace{4pt}

\par\begin{minipage}{\linewidth}\raggedright
% mdparagraph:P042:begin
\textbf{F13.}
% mdparagraph:P042:end

% item: NEW-U05-F13
% key: KEY-NEW-U05-F13
% mdtable:T15:begin
\par\begingroup
\fontsize{10.5}{13.2}\selectfont
\setlength{\LTleft}{0pt}\setlength{\LTright}{\fill}
\noindent
\begin{tabular}{@{}>{\raggedright\arraybackslash}p{38.43mm}>{\raggedright\arraybackslash}p{38.43mm}>{\raggedright\arraybackslash}p{41.51mm}@{}}
\toprule
\textbf{Present 2sg} & \textbf{Present 3sg} & \textbf{Preterite 1pl} \\
\midrule
færst & færð & fōron \\
\bottomrule
\end{tabular}
\par\endgroup

% mdtable:T15:end
% mdparagraph:P043:begin
The present singular mutation is \textbf{a → æ}. Class VI past \textbf{ō} is long and is shared by the singular and plural past stems. [4]
% mdparagraph:P043:end

\end{minipage}\par\addvspace{4pt}

\par\begin{minipage}{\linewidth}\raggedright
% mdparagraph:P044:begin
\textbf{F14.}
% mdparagraph:P044:end

% item: NEW-U05-F14
% key: KEY-NEW-U05-F14
% mdtable:T16:begin
\par\begingroup
\fontsize{10.5}{13.2}\selectfont
\setlength{\LTleft}{0pt}\setlength{\LTright}{\fill}
\noindent
\begin{tabular}{@{}>{\raggedright\arraybackslash}p{28.32mm}>{\raggedright\arraybackslash}p{28.32mm}>{\raggedright\arraybackslash}p{28.32mm}>{\raggedright\arraybackslash}p{30.59mm}@{}}
\toprule
\textbf{Present 1sg} & \textbf{Present 2sg} & \textbf{Present 2pl} & \textbf{Preterite 2sg} \\
\midrule
slēa & sliehst & slēað & slōge \\
\bottomrule
\end{tabular}
\par\endgroup

% mdtable:T16:end
% mdparagraph:P045:begin
The contracted present first singular is \textbf{slēa}, and the plural is \textbf{slēað}. \textbf{Sliehst} has short \textbf{ie} and \textbf{h}; plural-based past \textbf{slōge} has long \textbf{ō} and \textbf{g}. The documented first/\allowbreak{}third-singular past alternative \textbf{slōh} does not replace the required \textbf{g} in this second singular. [4]
% mdparagraph:P045:end

\end{minipage}\par\addvspace{4pt}

\par\begin{minipage}{\linewidth}\raggedright
% mdparagraph:P046:begin
\textbf{F15.}
% mdparagraph:P046:end

% item: NEW-U05-F15
% key: KEY-NEW-U05-F15
% mdtable:T17:begin
\par\begingroup
\fontsize{10.5}{13.2}\selectfont
\setlength{\LTleft}{0pt}\setlength{\LTright}{\fill}
\noindent
\begin{tabular}{@{}>{\raggedright\arraybackslash}p{38.43mm}>{\raggedright\arraybackslash}p{38.43mm}>{\raggedright\arraybackslash}p{41.51mm}@{}}
\toprule
\textbf{Present 2sg} & \textbf{Present 3sg} & \textbf{Preterite 3sg} \\
\midrule
stęntst & stęnt & stōd \\
\bottomrule
\end{tabular}
\par\endgroup

% mdtable:T17:end
% mdparagraph:P047:begin
The present singular has mutated \textbf{ę} and the dental combinations \textbf{-tst/\allowbreak{}-t}. The past loses \textbf{n} and has long \textbf{ō}: \textbf{stōd}. Adding a weak dental suffix to \textbf{stand-} would not produce this strong preterite. [4]
% mdparagraph:P047:end

\end{minipage}\par\addvspace{4pt}

\par\begin{minipage}{\linewidth}\raggedright
% mdparagraph:P048:begin
\textbf{F16.}
% mdparagraph:P048:end

% item: NEW-U05-F16
% key: KEY-NEW-U05-F16
% mdtable:T18:begin
\par\begingroup
\fontsize{10.5}{13.2}\selectfont
\setlength{\LTleft}{0pt}\setlength{\LTright}{\fill}
\noindent
\begin{tabular}{@{}>{\raggedright\arraybackslash}p{38.43mm}>{\raggedright\arraybackslash}p{38.43mm}>{\raggedright\arraybackslash}p{41.51mm}@{}}
\toprule
\textbf{Present 1sg} & \textbf{Present 3sg} & \textbf{Preterite 1pl} \\
\midrule
hębbe & hęfð & hōfon \\
\bottomrule
\end{tabular}
\par\endgroup

% mdtable:T18:end
% mdparagraph:P049:begin
The present alternates \textbf{bb} with \textbf{f} before the selected third-singular ending. The past uses \textbf{hōf-}, not the present vowel or doubled consonant. [4]
% mdparagraph:P049:end

\end{minipage}\par\addvspace{4pt}

\par\begin{minipage}{\linewidth}\raggedright
% mdparagraph:P050:begin
\textbf{F17.}
% mdparagraph:P050:end

% item: NEW-U05-F17
% key: KEY-NEW-U05-F17
% mdtable:T19:begin
\par\begingroup
\fontsize{10.5}{13.2}\selectfont
\setlength{\LTleft}{0pt}\setlength{\LTright}{\fill}
\noindent
\begin{tabular}{@{}>{\raggedright\arraybackslash}p{21.85mm}>{\raggedright\arraybackslash}p{21.85mm}>{\raggedright\arraybackslash}p{21.85mm}>{\raggedright\arraybackslash}p{23.6mm}>{\raggedright\arraybackslash}p{23.6mm}@{}}
\toprule
\textbf{Verb} & \textbf{Present 1sg} & \textbf{Present 2sg} & \textbf{Preterite 3sg} & \textbf{Preterite 3pl} \\
\midrule
scieppan & scieppe & sciepst & scēop & — \\
gescieppan & — & — & — & gescēopon \\
\bottomrule
\end{tabular}
\par\endgroup

% mdtable:T19:end
% mdparagraph:P051:begin
The doubled \textbf{pp} remains before the vowel ending in \textbf{scieppe} but is reduced before \textbf{-st}. Selected past \textbf{ēo} is long. Cook also records \textbf{scōp/\allowbreak{}scōpon}; the corresponding prefixed \textbf{gescōp/\allowbreak{}gescōpon} series is a documented grammatical expansion. These are valid alternatives, while the task selects the \textbf{ēo} series. The dashes mark cells not requested here. [4]
% mdparagraph:P051:end

\end{minipage}\par\addvspace{4pt}

\par\begin{minipage}{\linewidth}\raggedright
% mdparagraph:P052:begin
\textbf{F18.}
% mdparagraph:P052:end

% item: NEW-U05-F18
% key: KEY-NEW-U05-F18
% mdtable:T20:begin
\par\begingroup
\fontsize{10.5}{13.2}\selectfont
\setlength{\LTleft}{0pt}\setlength{\LTright}{\fill}
\noindent
\begin{tabular}{@{}>{\raggedright\arraybackslash}p{28.32mm}>{\raggedright\arraybackslash}p{28.32mm}>{\raggedright\arraybackslash}p{28.32mm}>{\raggedright\arraybackslash}p{30.59mm}@{}}
\toprule
\textbf{Present 1sg} & \textbf{Present 2sg} & \textbf{Present 3sg} & \textbf{Preterite 2sg} \\
\midrule
swęrie & swęrest & swęreð & swōre \\
\bottomrule
\end{tabular}
\par\endgroup

% mdtable:T20:end
% mdparagraph:P053:begin
The present follows the explicitly taught \textbf{-ie/\allowbreak{}-est/\allowbreak{}-eð} pattern; it is not made by keeping \textbf{-i-} before every ending. The preterite is strong, with \textbf{swōr-} and second-singular \textbf{-e}, not a weak dental past. [4]
% mdparagraph:P053:end

\end{minipage}\par\addvspace{4pt}

\par\begin{minipage}{\linewidth}\raggedright
% mdparagraph:P054:begin
\textbf{F19.}
% mdparagraph:P054:end

% item: NEW-U05-F19
% key: KEY-NEW-U05-F19
% mdtable:T21:begin
\par\begingroup
\fontsize{10.5}{13.2}\selectfont
\setlength{\LTleft}{0pt}\setlength{\LTright}{\fill}
\noindent
\begin{tabular}{@{}>{\raggedright\arraybackslash}p{38.43mm}>{\raggedright\arraybackslash}p{38.43mm}>{\raggedright\arraybackslash}p{41.51mm}@{}}
\toprule
\textbf{Verb} & \textbf{Present 3sg} & \textbf{Preterite 1pl} \\
\midrule
feallan & fielð & — \\
healdan & hielt & — \\
weaxan & wiext & wēoxon \\
\bottomrule
\end{tabular}
\par\endgroup

% mdtable:T21:end
% mdparagraph:P055:begin
All three present forms have short \textbf{ie} from short \textbf{ea}. \textbf{Fielð} reduces doubled \textbf{ll}; \textbf{hielt} has the dental outcome; \textbf{wiext} combines stem \textbf{x} with the ending as \textbf{xt}. The past \textbf{ēo} of \textbf{wēoxon} is long. [5]
% mdparagraph:P055:end

\end{minipage}\par\addvspace{4pt}

\par\begin{minipage}{\linewidth}\raggedright
% mdparagraph:P056:begin
\textbf{F20.}
% mdparagraph:P056:end

% item: NEW-U05-F20
% key: KEY-NEW-U05-F20
% mdtable:T22:begin
\par\begingroup
\fontsize{10.5}{13.2}\selectfont
\setlength{\LTleft}{0pt}\setlength{\LTright}{\fill}
\noindent
\begin{tabular}{@{}>{\raggedright\arraybackslash}p{38.43mm}>{\raggedright\arraybackslash}p{38.43mm}>{\raggedright\arraybackslash}p{41.51mm}@{}}
\toprule
\textbf{Verb} & \textbf{Present 2sg} & \textbf{Preterite 2sg} \\
\midrule
cnāwan & cnǣwst & cnēowe \\
wēpan & wēpst & wēope \\
\bottomrule
\end{tabular}
\par\endgroup

% mdtable:T22:end
% mdparagraph:P057:begin
\textbf{Cnāwan} changes present \textbf{ā} to \textbf{ǣ} in the requested present singular. \textbf{Wēpan} already has \textbf{ē} in the present; do not manufacture an additional changed vowel there. Both requested preterites use \textbf{ēo} plus the second-singular \textbf{-e}. [5]
% mdparagraph:P057:end

\end{minipage}\par\addvspace{4pt}

\par\begin{minipage}{\linewidth}\raggedright
% mdparagraph:P058:begin
\textbf{F21.}
% mdparagraph:P058:end

% item: NEW-U05-F21
% key: KEY-NEW-U05-F21
% mdtable:T23:begin
\par\begingroup
\fontsize{10.5}{13.2}\selectfont
\setlength{\LTleft}{0pt}\setlength{\LTright}{\fill}
\noindent
\begin{tabular}{@{}>{\raggedright\arraybackslash}p{38.43mm}>{\raggedright\arraybackslash}p{38.43mm}>{\raggedright\arraybackslash}p{41.51mm}@{}}
\toprule
\textbf{Verb} & \textbf{Present 2sg} & \textbf{Preterite 2sg} \\
\midrule
hēawan & hīewst & hēowe \\
flōwan & flēwst & flēowe \\
\bottomrule
\end{tabular}
\par\endgroup

% mdtable:T23:end
% mdparagraph:P059:begin
Long present \textbf{ēa → īe} in \textbf{hīewst} contrasts with \textbf{ō → ē} in \textbf{flēwst}. Both past forms have long \textbf{ēo}. Shortening \textbf{īe} in \textbf{hīewst} would erase the quantity distinction taught with these present changes. [5]
% mdparagraph:P059:end

\end{minipage}\par\addvspace{4pt}

\par\begin{minipage}{\linewidth}\raggedright
% mdparagraph:P060:begin
\textbf{F22.}
% mdparagraph:P060:end

% item: NEW-U05-F22
% key: KEY-NEW-U05-F22
% mdtable:T24:begin
\par\begingroup
\fontsize{10.5}{13.2}\selectfont
\setlength{\LTleft}{0pt}\setlength{\LTright}{\fill}
\noindent
\begin{tabular}{@{}>{\raggedright\arraybackslash}p{38.43mm}>{\raggedright\arraybackslash}p{38.43mm}>{\raggedright\arraybackslash}p{41.51mm}@{}}
\toprule
\textbf{Verb} & \textbf{Present 3sg} & \textbf{Preterite 2sg} \\
\midrule
lǣtan & lǣtt & — \\
hātan & hǣtt & — \\
fōn & fēhð & fēnge \\
\bottomrule
\end{tabular}
\par\endgroup

% mdtable:T24:end
% mdparagraph:P061:begin
Selected \textbf{lǣtt/\allowbreak{}hǣtt} have the double dental; \textbf{lǣt/\allowbreak{}hǣt} are documented alternatives. \textbf{Hǣtt} belongs to the active indicative and is not the separate \textbf{hātte} system. Contracted \textbf{fōn} has present \textbf{h} and plural-based past \textbf{ng}: \textbf{fēnge}. [5]
% mdparagraph:P061:end

\end{minipage}\par\addvspace{4pt}

\par\begin{minipage}{\linewidth}\raggedright
% mdparagraph:P062:begin
\textbf{F23.}
% mdparagraph:P062:end

% item: NEW-U05-F23
% key: KEY-NEW-U05-F23
% mdtable:T25:begin
\par\begingroup
\fontsize{10.5}{13.2}\selectfont
\setlength{\LTleft}{0pt}\setlength{\LTright}{\fill}
\noindent
\begin{tabular}{@{}>{\raggedright\arraybackslash}p{38.43mm}>{\raggedright\arraybackslash}p{38.43mm}>{\raggedright\arraybackslash}p{41.51mm}@{}}
\toprule
\textbf{Verb} & \textbf{Present 3sg} & \textbf{Preterite 2sg} \\
\midrule
geswīcan & geswīcð & geswice \\
scīnan & scīnð & scine \\
\bottomrule
\end{tabular}
\par\endgroup

% mdtable:T25:end
% mdparagraph:P063:begin
These Class I verbs keep long \textbf{ī} in the present but use short \textbf{i} in the plural-based past. \textbf{Geswāc/\allowbreak{}scān} are first/\allowbreak{}third singular past forms, not second singular. Learning an earlier class pattern did not make these new lexical forms previously taught. [6]
% mdparagraph:P063:end

\end{minipage}\par\addvspace{4pt}

\par\begin{minipage}{\linewidth}\raggedright
% mdparagraph:P064:begin
\textbf{F24.}
% mdparagraph:P064:end

% item: NEW-U05-F24
% key: KEY-NEW-U05-F24
% mdtable:T26:begin
\par\begingroup
\fontsize{10.5}{13.2}\selectfont
\setlength{\LTleft}{0pt}\setlength{\LTright}{\fill}
\noindent
\begin{tabular}{@{}>{\raggedright\arraybackslash}p{38.43mm}>{\raggedright\arraybackslash}p{38.43mm}>{\raggedright\arraybackslash}p{41.51mm}@{}}
\toprule
\textbf{Verb} & \textbf{Present 3sg} & \textbf{Preterite 3pl} \\
\midrule
flēogan & flīehð & flugon \\
swimman & swimð & swummon \\
\bottomrule
\end{tabular}
\par\endgroup

% mdtable:T26:end
% mdparagraph:P065:begin
Class II \textbf{flēogan} has long \textbf{īe} in the selected present singular and short \textbf{u} in the plural past. \textbf{Flīegð} is a documented present alternative. Class III \textbf{swimman} reduces \textbf{mm} before the present consonantal ending, but preserves it in plural past \textbf{swummon}. The documented \textbf{swommon} alternative differs from the selected \textbf{a/\allowbreak{}u} past series. [6]
% mdparagraph:P065:end

\end{minipage}\par\addvspace{4pt}

\reserve{7\baselineskip}
% mdparagraph:P066:begin
\section*{Clauses}
% mdparagraph:P066:end

\reserve{5\baselineskip}
\par\begin{minipage}{\linewidth}\raggedright
% mdparagraph:P067:begin
\textbf{C1.}
% mdparagraph:P067:end

% item: NEW-U05-C01
% key: KEY-NEW-U05-C01
% mdparagraph:P068:begin
\textbf{Ðū birst ðone męte, and ic nime ðone fisc.}
% mdparagraph:P068:end

% mdparagraph:P069:begin
The first subject is one addressee, with second-singular \textbf{birst}; the second is the speaker, with first-singular \textbf{nime}. Food and fish are accusative masculine singular. \textbf{Bierst} is a documented alternative to \textbf{birst}. Plural verbs would change the requested subjects, and a dative object would miss the taught carry/\allowbreak{}take constructions. [7]
% mdparagraph:P069:end

\end{minipage}\par\addvspace{4pt}

\par\begin{minipage}{\linewidth}\raggedright
% mdparagraph:P070:begin
\textbf{C2.}
% mdparagraph:P070:end

% item: NEW-U05-C02
% key: KEY-NEW-U05-C02
% mdparagraph:P071:begin
\textbf{Wit gēafon ðǣm mannum uncerne męte.}
% mdparagraph:P071:end

% mdparagraph:P072:begin
\textbf{Wit} and \textbf{uncerne} both preserve “us two.” The dual subject takes plural \textbf{gēafon}. \textbf{Ðǣm mannum} is the dative plural recipient; \textbf{uncerne męte} is the accusative masculine singular food given. The possessive ending agrees with \textbf{męte}, while its stem identifies the two possessors. Ordinary \textbf{wē} or \textbf{ūre} would lose the specified dual meaning. A grammatical sentence with \textbf{sęllan} would also miss the explicit \textbf{giefan} target. [7]
% mdparagraph:P072:end

\end{minipage}\par\addvspace{4pt}

\par\begin{minipage}{\linewidth}\raggedright
% mdparagraph:P073:begin
\textbf{C3.}
% mdparagraph:P073:end

% item: NEW-U05-C03
% key: KEY-NEW-U05-C03
% mdparagraph:P074:begin
\textbf{Ðā gefēran sǣton on ðǣm hofe, and se mann stōd on ðǣm strande.}
% mdparagraph:P074:end

% mdparagraph:P075:begin
The companions are plural and require \textbf{sǣton}; the man is singular and requires \textbf{stōd}. Both phrases with \textbf{on} locate the posture and take the dative. They do not describe movement into the dwelling or toward the shore. \textbf{Stōd} here means ordinary “stood”; the earlier reading's opposition sense is not needed. [7]
% mdparagraph:P075:end

\end{minipage}\par\addvspace{4pt}

\par\begin{minipage}{\linewidth}\raggedright
% mdparagraph:P076:begin
\textbf{C4.}
% mdparagraph:P076:end

% item: NEW-U05-C04
% key: KEY-NEW-U05-C04
% mdparagraph:P077:begin
\textbf{Ic ne ǣt ðone męte; ðū ǣte hine.}
% mdparagraph:P077:end

% mdparagraph:P078:begin
First-singular \textbf{ǣt} contrasts with second-singular \textbf{ǣte}. Only the first clause is negated. \textbf{Hine} is accusative masculine singular, referring to masculine \textbf{męte}; English “it” does not automatically require Old English neuter \textbf{hit}. Repeating \textbf{ðone męte} can preserve the reference, but the pronoun model demonstrates its agreement. Short singular \textbf{æt} is a documented alternative to selected \textbf{ǣt} in the first clause. [3, 7]
% mdparagraph:P078:end

\end{minipage}\par\addvspace{4pt}

\par\begin{minipage}{\linewidth}\raggedright
% mdparagraph:P079:begin
\textbf{C5.}
% mdparagraph:P079:end

% item: NEW-U05-C05
% key: KEY-NEW-U05-C05
% mdparagraph:P080:begin
\textbf{Hīe gesāwon ðā fuglas and nōmon ðā fiscas.}
% mdparagraph:P080:end

% mdparagraph:P081:begin
The plural subject controls both plural preterites. Both objects also become plural: \textbf{fuglas} and \textbf{fiscas}, each with accusative plural \textbf{ðā}. Changing only the verbs, or only the nouns, leaves part of the requested transformation undone. The model retains the original \textbf{gesēon} and the selected \textbf{nōm/\allowbreak{}nōmon} series; documented \textbf{nāmon} is a valid alternative past of \textbf{niman}.
% mdparagraph:P081:end

\end{minipage}\par\addvspace{4pt}

\par\begin{minipage}{\linewidth}\raggedright
% mdparagraph:P082:begin
\textbf{C6.}
% mdparagraph:P082:end

% item: NEW-U05-C06
% key: KEY-NEW-U05-C06
% mdparagraph:P083:begin
\textbf{Se mann cwið: “Ic sēo ðā trēowu.”}
% mdparagraph:P083:end

% mdparagraph:P084:begin
\textbf{Cwið} is third-singular present, governed by \textbf{se mann}. Inside the quotation, \textbf{ic} is the man speaking and \textbf{sēo} is first-singular present; \textbf{ðā trēowu} is accusative plural. \textbf{Gesēo} is an equally acceptable separately taught seeing verb here. A past “said/\allowbreak{}saw” changes the specified tense. The task requires direct speech introduced by \textbf{cweðan}, not a newly formed indirect-speech construction. [7]
% mdparagraph:P084:end

\end{minipage}\par\addvspace{4pt}

\par\begin{minipage}{\linewidth}\raggedright
% mdparagraph:P085:begin
\textbf{C7.}
% mdparagraph:P085:end

% item: NEW-U05-C07
% key: KEY-NEW-U05-C07
% mdparagraph:P086:begin
\textbf{Gē geswīcað ēowres weorces.}
% mdparagraph:P086:end

% mdparagraph:P087:begin
\textbf{Gē} and \textbf{geswīcað} are second-person plural. \textbf{Geswīcan} governs the genitive of the activity ceased from: \textbf{weorces}. The agreeing possessive \textbf{ēowres} is genitive neuter singular and identifies plural addressees as possessors. English “from” is expressed by this genitive construction, without adding \textbf{fram}. [8]
% mdparagraph:P087:end

\end{minipage}\par\addvspace{4pt}

\par\begin{minipage}{\linewidth}\raggedright
% mdparagraph:P088:begin
\textbf{C8.}
% mdparagraph:P088:end

% item: NEW-U05-C08
% key: KEY-NEW-U05-C08
% mdparagraph:P089:begin
\textbf{Ðæt trēo fielð, and ðā męnn healdað ðæt scip.}
% mdparagraph:P089:end

% mdparagraph:P090:begin
The singular tree is the subject of intransitive \textbf{fielð}. The plural men are the subject of \textbf{healdað}, and the singular ship is its accusative object. Although \textbf{scip} has identical nominative and accusative singular spelling, its role here is determined by the clause. A plural ship phrase or a singular holding verb would change the requested relations. [7]
% mdparagraph:P090:end

\end{minipage}\par\addvspace{4pt}

\reserve{7\baselineskip}
% mdparagraph:P091:begin
\section*{Cook, Group I, Exercise 5}
% mdparagraph:P091:end

% mdparagraph:P092:begin
The English prompts below retain Cook's wording and original item numbers, printed p.3 /\allowbreak{} PDF11. The common target is a complete translation with the correct finite tense and agreement, noun and adjective cases, and stated lexical relations. The models use the selected forms; the notes distinguish any narrower taught target from other valid language. The source has no attached item notes. [9]
% mdparagraph:P092:end

\reserve{5\baselineskip}
\par\begin{minipage}{\linewidth}\raggedright
% mdparagraph:P093:begin
\textbf{I.5.1.}
% mdparagraph:P093:end

% item: EX-I-05-01
% key: KEY-EX-I-05-01
% mdparagraph:P094:begin
God created the heavens and the earth.
% mdparagraph:P094:end

% mdparagraph:P095:begin
\textbf{God gescēop ðā heofonas and ðā eorðan.}
% mdparagraph:P095:end

% mdparagraph:P096:begin
\textbf{God} governs singular preterite \textbf{gescēop}. \textbf{Heofonas} is accusative masculine plural; \textbf{eorðan} is accusative feminine singular. Their articles happen to share the spelling \textbf{ðā}. \textbf{Scēop} is also acceptable, and the documented \textbf{gescōp/\allowbreak{}scōp} alternatives retain the creation meaning. The reading's weak-feminine singular \textbf{heofonan} does not preserve this English prompt's plural “heavens.”
% mdparagraph:P096:end

\end{minipage}\par\addvspace{4pt}

\par\begin{minipage}{\linewidth}\raggedright
% mdparagraph:P097:begin
\textbf{I.5.2.}
% mdparagraph:P097:end

% item: EX-I-05-02
% key: KEY-EX-I-05-02
% mdparagraph:P098:begin
God made two great lights.
% mdparagraph:P098:end

% mdparagraph:P099:begin
\textbf{God geworhte twā miclu lēoht.}
% mdparagraph:P099:end

% mdparagraph:P100:begin
The object is accusative neuter plural: \textbf{twā}, strong \textbf{miclu}, and \textbf{lēoht}. \textbf{Lēoht} does not gain a plural \textbf{-as} or \textbf{-u} in its selected long-stem neuter paradigm. The numeral does not by itself require a weak adjective. \textbf{Geworhte} is the already taught weak preterite of \textbf{gewyrcean}.
% mdparagraph:P100:end

\end{minipage}\par\addvspace{4pt}

\par\begin{minipage}{\linewidth}\raggedright
% mdparagraph:P101:begin
\textbf{I.5.3.}
% mdparagraph:P101:end

% item: EX-I-05-03
% key: KEY-EX-I-05-03
% mdparagraph:P102:begin
God made the firmament.
% mdparagraph:P102:end

% mdparagraph:P103:begin
\textbf{God geworhte ðā fæstnesse.}
% mdparagraph:P103:end

% mdparagraph:P104:begin
\textbf{Ðā fæstnesse} is accusative feminine singular, the thing made. The finite verb is third-singular preterite. Its dental past formation remains weak even though this unit chiefly introduces strong systems.
% mdparagraph:P104:end

\end{minipage}\par\addvspace{4pt}

\par\begin{minipage}{\linewidth}\raggedright
% mdparagraph:P105:begin
\textbf{I.5.4.}
% mdparagraph:P105:end

% item: EX-I-05-04
% key: KEY-EX-I-05-04
% mdparagraph:P106:begin
God created the great whales.
% mdparagraph:P106:end

% mdparagraph:P107:begin
\textbf{God gescēop ðā miclan hwalas.}
% mdparagraph:P107:end

% mdparagraph:P108:begin
The object is accusative masculine plural. \textbf{Miclan} is weak after \textbf{ðā}, and \textbf{hwalas} has the selected plural stem \textbf{hwal-}, not singular \textbf{hwæl-}. \textbf{Scēop} and the documented \textbf{ō}-series creation forms are valid alternatives in this unqualified English prompt. A plural verb would incorrectly agree with the object instead of \textbf{God}.
% mdparagraph:P108:end

\end{minipage}\par\addvspace{4pt}

\par\begin{minipage}{\linewidth}\raggedright
% mdparagraph:P109:begin
\textbf{I.5.5.}
% mdparagraph:P109:end

% item: EX-I-05-05
% key: KEY-EX-I-05-05
% mdparagraph:P110:begin
He made the beasts of the earth.
% mdparagraph:P110:end

% mdparagraph:P111:begin
\textbf{Hē geworhte ðā dēor ðǣre eorðan.}
% mdparagraph:P111:end

% mdparagraph:P112:begin
\textbf{Ðā dēor} is accusative neuter plural. \textbf{Ðǣre eorðan} is the dependent genitive singular, “of the earth.” The spelling \textbf{dēor} is the same in singular and plural, so the article and intended meaning matter. The limiting genitive may instead precede its head: \textbf{Hē geworhte ðǣre eorðan dēor.} Both arrangements preserve the relation.
% mdparagraph:P112:end

\end{minipage}\par\addvspace{4pt}

\par\begin{minipage}{\linewidth}\raggedright
% mdparagraph:P113:begin
\textbf{I.5.6.}
% mdparagraph:P113:end

% item: EX-I-05-06
% key: KEY-EX-I-05-06
% mdparagraph:P114:begin
He saw all those things.
% mdparagraph:P114:end

% mdparagraph:P115:begin
\textbf{Hē geseah eall ðā ðing.}
% mdparagraph:P115:end

% mdparagraph:P116:begin
The independent target is the subject and third-singular preterite seeing form. \textbf{Seah} is also acceptable. The complete accusative plural phrase \textbf{eall ðā ðing} is supplied and may be copied; do not require the learner to reconstruct its quantifier syntax or endings from memory. Present \textbf{gesēo/\allowbreak{}gesiehð} would miss “saw.”
% mdparagraph:P116:end

\end{minipage}\par\addvspace{4pt}

\par\begin{minipage}{\linewidth}\raggedright
% mdparagraph:P117:begin
\textbf{I.5.7.}
% mdparagraph:P117:end

% item: EX-I-05-07
% key: KEY-EX-I-05-07
% mdparagraph:P118:begin
He finished his work.
% mdparagraph:P118:end

% mdparagraph:P119:begin
\textbf{Hē gefylde his weorc.}
% mdparagraph:P119:end

% mdparagraph:P120:begin
\textbf{Gefylde} is third-singular weak preterite of \textbf{gefyllan}, in its taught finish/\allowbreak{}accomplish sense. \textbf{Weorc} is accusative neuter singular. \textbf{His} identifies a male possessor and does not acquire a neuter accusative ending; the natural reference is the subject, though a context could identify another male owner.
% mdparagraph:P120:end

\end{minipage}\par\addvspace{4pt}

\par\begin{minipage}{\linewidth}\raggedright
% mdparagraph:P121:begin
\textbf{I.5.8.}
% mdparagraph:P121:end

% item: EX-I-05-08
% key: KEY-EX-I-05-08
% mdparagraph:P122:begin
He blesses the seventh day.
% mdparagraph:P122:end

% mdparagraph:P123:begin
\textbf{Hē geblētsað ðone seofoðan dæg.}
% mdparagraph:P123:end

% mdparagraph:P124:begin
\textbf{Geblētsað} is the taught weak Class II third-singular present; its \textbf{-að} does not make it plural in this system. \textbf{Ðone seofoðan dæg} is accusative masculine singular, with the weak ordinal after the article. A preterite would change “blesses” to “blessed.”
% mdparagraph:P124:end

\end{minipage}\par\addvspace{4pt}

\par\begin{minipage}{\linewidth}\raggedright
% mdparagraph:P125:begin
\textbf{I.5.9.}
% mdparagraph:P125:end

% item: EX-I-05-09
% key: KEY-EX-I-05-09
% mdparagraph:P126:begin
The man ceases from his work.
% mdparagraph:P126:end

% mdparagraph:P127:begin
\textbf{Se mann geswīcð his weorces.}
% mdparagraph:P127:end

% mdparagraph:P128:begin
\textbf{Geswīcð} is third-singular present. \textbf{Weorces} is genitive under \textbf{geswīcan}, and \textbf{his} remains invariant. Do not substitute the past \textbf{geswāc}, or add a preposition merely because English has “from.” [8]
% mdparagraph:P128:end

\end{minipage}\par\addvspace{4pt}

\par\begin{minipage}{\linewidth}\raggedright
% mdparagraph:P129:begin
\textbf{I.5.10.}
% mdparagraph:P129:end

% item: EX-I-05-10
% key: KEY-EX-I-05-10
% mdparagraph:P130:begin
I see the bird in the tree.
% mdparagraph:P130:end

% mdparagraph:P131:begin
\textbf{Ic gesēo ðone fugol on ðǣm trēowe.}
% mdparagraph:P131:end

% mdparagraph:P132:begin
\textbf{Sēo} is also acceptable for the first-singular present. The bird is accusative masculine singular; the location phrase has dative neuter singular \textbf{ðǣm trēowe}. The noun's selected dative retains \textbf{w} and \textbf{-e}. There is no motion into the tree.
% mdparagraph:P132:end

\end{minipage}\par\addvspace{4pt}

\par\begin{minipage}{\linewidth}\raggedright
% mdparagraph:P133:begin
\textbf{I.5.11.}
% mdparagraph:P133:end

% item: EX-I-05-11
% key: KEY-EX-I-05-11
% mdparagraph:P134:begin
They saw the fish in the water.
% mdparagraph:P134:end

% mdparagraph:P135:begin
Two equally valid number choices are:
% mdparagraph:P135:end

% mdtable:T27:begin
\par\begingroup
\fontsize{10.5}{13.2}\selectfont
\setlength{\LTleft}{0pt}\setlength{\LTright}{\fill}
\noindent
\begin{tabular}{@{}>{\raggedright\arraybackslash}p{48.98mm}>{\raggedright\arraybackslash}p{72.21mm}@{}}
\toprule
\textbf{Fish understood as} & \textbf{Model} \\
\midrule
Singular & Hīe gesāwon ðone fisc on ðǣm wætere. \\
Plural & Hīe gesāwon ðā fiscas on ðǣm wætere. \\
\bottomrule
\end{tabular}
\par\endgroup

% mdtable:T27:end
% mdparagraph:P136:begin
English \textbf{fish} does not settle the object's number, and Cook supplies no note restricting it. Both \textbf{ðone fisc} and \textbf{ðā fiscas} meet the prompt. \textbf{Sāwon} is also acceptable in either model. The subject \textbf{they} requires plural \textbf{hīe} and a plural preterite; object-number freedom does not permit singular \textbf{geseah/\allowbreak{}seah}. The water phrase is dative location.
% mdparagraph:P136:end

\end{minipage}\par\addvspace{4pt}

\par\begin{minipage}{\linewidth}\raggedright
% mdparagraph:P137:begin
\textbf{I.5.12.}
% mdparagraph:P137:end

% item: EX-I-05-12
% key: KEY-EX-I-05-12
% mdparagraph:P138:begin
The birds of the air fly.
% mdparagraph:P138:end

% mdparagraph:P139:begin
\textbf{Ðæs lyftes fuglas flēogað.}
% mdparagraph:P139:end

% mdparagraph:P140:begin
\textbf{Fuglas} is nominative plural and controls present \textbf{flēogað}. \textbf{Ðæs lyftes} is genitive singular of the selected neuter \textbf{lyft}, identifying the birds as “of the air.” \textbf{Ðā fuglas ðæs lyftes flēogað} also preserves the relation. The genitive is not an object of “fly”; this is ordinary intransitive winged flying. [7]
% mdparagraph:P140:end

\end{minipage}\par\addvspace{4pt}

\par\begin{minipage}{\linewidth}\raggedright
% mdparagraph:P141:begin
\textbf{I.5.13.}
% mdparagraph:P141:end

% item: EX-I-05-13
% key: KEY-EX-I-05-13
% mdparagraph:P142:begin
The trees grew.
% mdparagraph:P142:end

% mdparagraph:P143:begin
\textbf{Ðā trēowu wēoxon.}
% mdparagraph:P143:end

% mdparagraph:P144:begin
The trees are nominative neuter plural, with the selected \textbf{trēowu}. \textbf{Wēoxon} is third-plural preterite of the taught growth verb \textbf{weaxan}, with long \textbf{ēo}. Present \textbf{weaxað} or singular past \textbf{wēox} misses a requested distinction. Cook also lists \textbf{grōwan} and explicitly assigns it the \textbf{flōwan} pattern, giving the valid alternative \textbf{grēowon} by that pattern. Its full system is not an independent-recall requirement here. It differs from the selected \textbf{weaxan} practice target, not from the English meaning. [5, 9]
% mdparagraph:P144:end

\end{minipage}\par\addvspace{4pt}

\par\begin{minipage}{\linewidth}\raggedright
% mdparagraph:P145:begin
\textbf{I.5.14.}
% mdparagraph:P145:end

% item: EX-I-05-14
% key: KEY-EX-I-05-14
% mdparagraph:P146:begin
The stars shone.
% mdparagraph:P146:end

% mdparagraph:P147:begin
\textbf{Ðā steorran scinon.}
% mdparagraph:P147:end

% mdparagraph:P148:begin
Weak masculine \textbf{steorran} is nominative plural here, and \textbf{scinon} is third-plural preterite with short \textbf{i}. Present \textbf{scīnað} has long \textbf{ī}; singular past \textbf{scān} has long \textbf{ā}. The target is literal shining, not a new figurative sense.
% mdparagraph:P148:end

\end{minipage}\par\addvspace{4pt}

\par\begin{minipage}{\linewidth}\raggedright
% mdparagraph:P149:begin
\textbf{I.5.15.}
% mdparagraph:P149:end

% item: EX-I-05-15
% key: KEY-EX-I-05-15
% mdparagraph:P150:begin
The fish swims.
% mdparagraph:P150:end

% mdparagraph:P151:begin
\textbf{Se fisc swimð.}
% mdparagraph:P151:end

% mdparagraph:P152:begin
Here English \textbf{swims} fixes the subject as singular. \textbf{Se fisc} is nominative masculine singular and controls third-singular \textbf{swimð}, with reduced \textbf{m} before the ending. The ambiguity allowed for the object in I.5.11 does not apply here.
% mdparagraph:P152:end

\end{minipage}\par\addvspace{4pt}

\par\begin{minipage}{\linewidth}\raggedright
% mdparagraph:P153:begin
\textbf{I.5.16.}
% mdparagraph:P153:end

% item: EX-I-05-16
% key: KEY-EX-I-05-16
% mdparagraph:P154:begin
The men swam.
% mdparagraph:P154:end

% mdparagraph:P155:begin
\textbf{Ðā męnn swummon.}
% mdparagraph:P155:end

% mdparagraph:P156:begin
\textbf{Męnn} is the nominative plural of \textbf{mann}. The plural preterite \textbf{swummon} retains doubled \textbf{mm} and has \textbf{u}; singular \textbf{swamm} and present \textbf{swimmað} miss the target. Documented \textbf{swommon} is a valid alternative past form, distinguished from the lesson's selected \textbf{a/\allowbreak{}u} series. [6]
% mdparagraph:P156:end

\end{minipage}\par\addvspace{4pt}

\reserve{7\baselineskip}
% mdparagraph:P157:begin
\section*{Food for the companions and the birds}
% mdparagraph:P157:end

\reserve{5\baselineskip}
\par\begin{minipage}{\linewidth}\raggedright
% mdparagraph:P158:begin
\textbf{T1.}
% mdparagraph:P158:end

% item: NEW-U05-T01
% key: KEY-NEW-U05-T01
% mdparagraph:P159:begin
One man gave his food to the birds. The birds ate and flew.
% mdparagraph:P159:end

% mdparagraph:P160:begin
\textbf{Ðǣm fuglum} identifies the dative plural recipients of \textbf{geaf}. In the next sentence, nominative \textbf{ðā fuglas} is the subject of both \textbf{ǣton} and \textbf{flugon}. “Flew away” is a possible contextual inference, but the text only states that they flew; departure is not a required detail. Giving food to the companions belongs to the earlier action by the men.
% mdparagraph:P160:end

\end{minipage}\par\addvspace{4pt}

\par\begin{minipage}{\linewidth}\raggedright
% mdparagraph:P161:begin
\textbf{T2.}
% mdparagraph:P161:end

% item: NEW-U05-T02
% key: KEY-NEW-U05-T02
% mdtable:T28:begin
\par\begingroup
\fontsize{10.5}{13.2}\selectfont
\setlength{\LTleft}{0pt}\setlength{\LTright}{\fill}
\noindent
\begin{tabular}{@{}>{\raggedright\arraybackslash}p{28.32mm}>{\raggedright\arraybackslash}p{39.79mm}>{\raggedright\arraybackslash}p{50.27mm}@{}}
\toprule
\textbf{Phrase} & \textbf{Case and number} & \textbf{Function} \\
\midrule
ðǣm gefērum & Dative plural & Recipients: the companions \\
ðone męte & Accusative singular & Thing given: the food \\
ðā męnn & Nominative plural & Subject of gēafon: the men \\
\bottomrule
\end{tabular}
\par\endgroup

% mdtable:T28:end
% mdparagraph:P162:begin
\textbf{Gēafon} agrees with the plural men who give. The recipients' plurality does not control its agreement; they are dative. English “gave the companions the food” and “gave the food to the companions” express the same relations. [7]
% mdparagraph:P162:end

\end{minipage}\par\addvspace{4pt}

\par\begin{minipage}{\linewidth}\raggedright
% mdparagraph:P163:begin
\textbf{T3.}
% mdparagraph:P163:end

% item: NEW-U05-T03
% key: KEY-NEW-U05-T03
% mdparagraph:P164:begin
\textbf{Se mann sitt on ðǣm strande and gesiehð ðā fiscas. Ðā fiscas swimmað on ðǣm wætere.}
% mdparagraph:P164:end

% mdparagraph:P165:begin
Singular past \textbf{sæt/\allowbreak{}geseah} becomes present \textbf{sitt/\allowbreak{}gesiehð}; plural past \textbf{swummon} becomes \textbf{swimmað}. The man remains the subject of both sitting and seeing. The fish remain the accusative plural object in the first sentence and the nominative plural subject in the second. Both location phrases retain the dative. Documented \textbf{sit} is also valid for the first verb; the model retains the account's \textbf{gesēon}.
% mdparagraph:P165:end

\end{minipage}\par\addvspace{4pt}

\par\begin{minipage}{\linewidth}\raggedright
% mdparagraph:P166:begin
\textbf{T4.}
% mdparagraph:P166:end

% item: NEW-U05-T04
% key: KEY-NEW-U05-T04
% mdparagraph:P167:begin
Two men carried food. Their companions sat in the dwelling. The men gave the companions the food, and they ate. One man did not eat. He gave his food to the birds. The birds ate and flew. The man sat on the shore and saw the fish. The fish swam in the water.
% mdparagraph:P167:end

% mdparagraph:P168:begin
This translation is independently prepared from the complete new Old English account. \textbf{Hiera} refers naturally to the two men as possessors of the companions, and \textbf{his męte} to the later man's food or share. In \textbf{and hīe ǣton}, the companions are the natural intended eaters, but the pronoun's form does not uniquely distinguish them from the men. Retain “they,” or accept a coherent contextual interpretation; do not require a uniquely determined antecedent. \textbf{Ān mann} likewise does not explicitly state whether the man belongs to the first pair or the companions.
% mdparagraph:P168:end

% mdparagraph:P169:begin
The verbs report past actions; “were sitting” or “were swimming” may be natural English renderings in context. \textbf{Etan} is used without an expressed object in the eating clauses, a taught use. The food's kind and the birds' destination are not specified. Unlike the unqualified English \textbf{fish} in Cook I.5.11, the account's \textbf{ðā fiscas} is explicitly plural in both occurrences. Preserve that number and the contrast between the man who did not eat and the birds who did. [7]
% mdparagraph:P169:end

\end{minipage}\par\addvspace{4pt}

\reserve{7\baselineskip}
% mdparagraph:P170:begin
\section*{Creation: first-day revisit}
% mdparagraph:P170:end

\reserve{9\baselineskip}
\par\begin{minipage}{\linewidth}\raggedright
% mdparagraph:P171:begin
\textbf{R1.}
% mdparagraph:P171:end

% item: NEW-U05-R01
% key: KEY-NEW-U05-R01
% mdtable:T29:begin
\par\begingroup
\fontsize{10.5}{13.2}\selectfont
\setlength{\LTleft}{0pt}\setlength{\LTright}{\fill}
\noindent
\begin{tabular}{@{}>{\raggedright\arraybackslash}p{21.78mm}>{\raggedright\arraybackslash}p{21.78mm}>{\raggedright\arraybackslash}p{21.78mm}>{\raggedright\arraybackslash}p{50.21mm}@{}}
\toprule
\textbf{Source form} & \textbf{Lemma} & \textbf{Strong class} & \textbf{Tense, person and number in the passage} \\
\midrule
gescēop & gescieppan & VI & Preterite, third singular \\
cwæð & cweðan & V & Preterite, third singular \\
geseah & gesēon & V & Preterite, third singular \\
hēt & hātan & VII & Preterite, third singular \\
\bottomrule
\end{tabular}
\par\endgroup

% mdtable:T29:end
% mdparagraph:P172:begin
All four are indicative. The first three have \textbf{God} as their stated subject. \textbf{Hēt} continues the same divine subject, understood from the preceding clauses; the light is what God names, not the namer. The bare verb spellings can also serve first-person singular elsewhere, but the passage fixes the person here. The naming construction remains supplied. [10]
% mdparagraph:P172:end

\end{minipage}\par\addvspace{4pt}

\par\begin{minipage}{\linewidth}\raggedright
% mdparagraph:P173:begin
\textbf{R2.}
% mdparagraph:P173:end

% item: NEW-U05-R02
% key: KEY-NEW-U05-R02
% mdparagraph:P174:begin
The ordinary taught meaning of \textbf{weorðan} is “become.” In the supplied \textbf{wearð geworht}, it combines with the participle to express “was made” or “came to be made.” \textbf{Wearð} is the finite word: third-person singular preterite indicative of \textbf{weorðan}. \textbf{Geworht} is the supplied past participle of \textbf{gewyrcean}.
% mdparagraph:P174:end

% mdparagraph:P175:begin
Recognizing the finite form is learned knowledge; forming a passive independently is not required. Do not parse \textbf{geworht} as a second finite preterite or translate the phrase as an active “made” with light as the maker. [10]
% mdparagraph:P175:end

\end{minipage}\par\addvspace{4pt}

\par\begin{minipage}{\linewidth}\raggedright
% mdparagraph:P176:begin
\textbf{R3.}
% mdparagraph:P176:end

% item: NEW-U05-R03
% key: KEY-NEW-U05-R03
% mdtable:T30:begin
\par\begingroup
\fontsize{10.5}{13.2}\selectfont
\setlength{\LTleft}{0pt}\setlength{\LTright}{\fill}
\noindent
\begin{tabular}{@{}>{\raggedright\arraybackslash}p{40.24mm}>{\raggedright\arraybackslash}p{36.85mm}>{\raggedright\arraybackslash}p{41.29mm}@{}}
\toprule
\textbf{Expression} & \textbf{Noun and gender} & \textbf{Number and case here} \\
\midrule
Source heofonan & heofone, weak feminine & Accusative singular: heaven \\
Exercise heavens: heofonas & heofon, strong masculine & Accusative plural: heavens \\
\bottomrule
\end{tabular}
\par\endgroup

% mdtable:T30:end
% mdparagraph:P177:begin
The source's \textbf{-an} belongs to a weak feminine singular form; it is not evidence for a plural “heavens” in that clause. The exercise calls for the already taught masculine plural. Both nouns refer to heaven, but their gender, inflection and number must be kept distinct. The source note supplies \textbf{heofone}; no new complete productive paradigm of that noun is required. [10]
% mdparagraph:P177:end

\end{minipage}\par\addvspace{4pt}

\par\begin{minipage}{\linewidth}\raggedright
% mdparagraph:P178:begin
\textbf{R4.}
% mdparagraph:P178:end

% item: NEW-U05-R04
% key: KEY-NEW-U05-R04
% mdparagraph:P179:begin
In the beginning God created heaven and earth. Now the earth was empty and void; and darkness was over the surface of the deep; and God's spirit was borne over the waters. God then said, “Let light come to be”; and light was made. God then saw that it was good; and he separated the light from the darkness. And he called the light Day, and the darkness Night. Then evening and morning had come to be—one day.
% mdparagraph:P179:end

% mdparagraph:P180:begin
This translation is prepared from all 69 words of the unchanged first-day passage. \textbf{Heaven} preserves source singular \textbf{heofonan}. “Now” renders connective \textbf{sōðlice} in context, following Cook's \textbf{autem} note; “indeed” is also defensible if the connection remains clear. “The expanse of the deep” is a more literal alternative to “the surface of the deep.” \textbf{Wæs gefęred} is the supplied “was carried/\allowbreak{}borne”; natural contextual “moved over” can convey the scene without changing God's spirit into the waters.
% mdparagraph:P180:end

% mdparagraph:P181:begin
The quotation expresses the supplied wish, “Let there be light” or “Let light come to be.” The following \textbf{wearð geworht} says that light was made. \textbf{Hit} in the content clause refers to the light God saw to be good. Preserve the separation and both names; do not stop at the making of light. The closing formula may be rendered “there came to be evening and morning—one day.” \textbf{Ān} says “one,” not the ordinal “first.” The whole closing construction remains supported, so a unique subject/\allowbreak{}predicate analysis of it is not a grading requirement. [10]
% mdparagraph:P181:end

\end{minipage}\par\addvspace{4pt}

\reserve{7\baselineskip}
% mdparagraph:P182:begin
\section*{Creation: second day}
% mdparagraph:P182:end

\reserve{9\baselineskip}
\par\begin{minipage}{\linewidth}\raggedright
% mdparagraph:P183:begin
\textbf{R5.}
% mdparagraph:P183:end

% item: NEW-U05-R05
% key: KEY-NEW-U05-R05
% mdtable:T31:begin
\par\begingroup
\fontsize{10.5}{13.2}\selectfont
\setlength{\LTleft}{0pt}\setlength{\LTright}{\fill}
\noindent
\begin{tabular}{@{}>{\raggedright\arraybackslash}p{21.15mm}>{\raggedright\arraybackslash}p{21.15mm}>{\raggedright\arraybackslash}p{35.21mm}>{\raggedright\arraybackslash}p{38.06mm}@{}}
\toprule
\textbf{Source form} & \textbf{Lemma} & \textbf{Formation} & \textbf{Person, number and tense here} \\
\midrule
cwæð & cweðan & Strong V & Third singular preterite \\
hēt & hātan & Strong VII & Third singular preterite \\
geworhte & gewyrcean & Weak & Third singular preterite \\
tōtwǣmde & tōtwǣman & Weak; supplied source form & Third singular preterite \\
\bottomrule
\end{tabular}
\par\endgroup

% mdtable:T31:end
% mdparagraph:P184:begin
All are indicative and share \textbf{God} as their subject. In coordinated \textbf{and tōtwǣmde}, God remains understood as the separator. \textbf{Cwæð/\allowbreak{}hēt} use strong past stems; \textbf{geworhte/\allowbreak{}tōtwǣmde} contain dental preterite formations. The \textbf{ð} in \textbf{cwæð} is part of its strong stem, not a weak preterite suffix. \textbf{Tōtwǣmde} may be identified from the local help; its full conjugation is not an independent production requirement. [10]
% mdparagraph:P184:end

\end{minipage}\par\addvspace{4pt}

\par\begin{minipage}{\linewidth}\raggedright
% mdparagraph:P185:begin
\textbf{R6.}
% mdparagraph:P185:end

% item: NEW-U05-R06
% key: KEY-NEW-U05-R06
% mdparagraph:P186:begin
God separates the waters \textbf{under the firmament} from those \textbf{above the firmament}. In \textbf{fram ðǣm ðe wǣron bufan ðǣre fæstnesse}, \textbf{ðǣm} is dative plural, “those,” with waters understood. It is used substantivally rather than as an article followed by an expressed noun.
% mdparagraph:P186:end

% mdparagraph:P187:begin
Both relative \textbf{ðe} occurrences are subjects of their own \textbf{wǣron} clauses: the respective waters “were under” and “were above.” They are not direct objects within those relative clauses. The surrounding verb has the first group as what is separated and \textbf{fram} introduces the group separated from. The local help supplies these constructions; independent formation of relative clauses is not tested. [10]
% mdparagraph:P187:end

\end{minipage}\par\addvspace{4pt}

\par\begin{minipage}{\linewidth}\raggedright
% mdparagraph:P188:begin
\textbf{R7.}
% mdparagraph:P188:end

% item: NEW-U05-R07
% key: KEY-NEW-U05-R07
% mdparagraph:P189:begin
God calls the firmament \textbf{Heaven}, source \textbf{Heofonan}. The supplied naming construction places both the named thing, \textbf{ðā fæstnesse}, and the name, \textbf{Heofonan}, in the accusative. \textbf{Heofonan} is again weak feminine singular, not a plural ending.
% mdparagraph:P189:end

% mdparagraph:P190:begin
In \textbf{ōðer dæg}, \textbf{ōðer} is masculine nominative singular, agreeing with masculine nominative singular \textbf{dæg}: “a second day.” The previously taught ordinal \textbf{ōðer} has this strong form. The larger evening-and-morning formula remains supplied; agreement within \textbf{ōðer dæg} does not require choosing one general clause analysis for the whole formula. [10]
% mdparagraph:P190:end

\end{minipage}\par\addvspace{4pt}

\par\begin{minipage}{\linewidth}\raggedright
% mdparagraph:P191:begin
\textbf{R8.}
% mdparagraph:P191:end

% item: NEW-U05-R08
% key: KEY-NEW-U05-R08
% mdparagraph:P192:begin
God then said again, “Let a firmament now come to be in the midst of the waters, and let it separate the waters from the waters.” And God made the firmament, and separated the waters which were under the firmament from those which were above the firmament; it was then done so. And God called the firmament Heaven. And then evening and morning had come to be—a second day.
% mdparagraph:P192:end

% mdparagraph:P193:begin
This translation is independently prepared from all 58 words of the second day. “God then also said” is a defensible alternative because Cook's note glosses \textbf{ęft} with \textbf{quoque}. Preserve \textbf{nū}, “now,” in the quotation. \textbf{Fæstnes}, the firmament, is the subject of both supplied wishes: it is to come into being and to separate the waters. “Let there now be a firmament” is an equally natural rendering of the first wish.
% mdparagraph:P193:end

% mdparagraph:P194:begin
The narrative then names God as the maker and separator. \textbf{Ðǣm} in the second relative phrase refers to the upper waters; it does not introduce another firmament. \textbf{Hit wæs ðā swā gedōn} reports fulfillment, “it was then done so,” or naturally “and so it was.” Keep it distinct from first-day \textbf{wearð geworht} and the final \textbf{wæs geworden} formula. The last clause says “a second day,” retaining \textbf{ōðer}. Both naming and the complete ending must be included, with the supplied help available throughout. [10]
% mdparagraph:P194:end

\end{minipage}\par\addvspace{4pt}

\clearpage
% mdparagraph:P195:begin
\section*{Cumulative review}
% mdparagraph:P195:end

\reserve{7\baselineskip}
\par\begin{minipage}{\linewidth}\raggedright
% mdparagraph:P196:begin
\textbf{X1.}
% mdparagraph:P196:end

% item: NEW-U05-X01
% key: KEY-NEW-U05-X01
% mdtable:T32:begin
\par\begingroup
\fontsize{10.5}{13.2}\selectfont
\setlength{\LTleft}{0pt}\setlength{\LTright}{\fill}
\noindent
\begin{tabular}{@{}>{\raggedright\arraybackslash}p{41.96mm}>{\raggedright\arraybackslash}p{79.22mm}@{}}
\toprule
\textbf{Third singular} & \textbf{Corresponding second singular} \\
\midrule
Hē geaf. & Ðū gēafe. \\
Hē sealde. & Ðū sealdest. \\
\bottomrule
\end{tabular}
\par\endgroup

% mdtable:T32:end
% mdparagraph:P197:begin
Both new clauses mean “You, one person, gave.” Strong \textbf{giefan} uses its plural-preterite stem \textbf{gēaf-} plus \textbf{-e}, including the long vowel. Weak \textbf{sęllan} keeps its dental preterite formation and adds second-singular \textbf{-est}. Applying the weak ending to \textbf{geaf}, or leaving \textbf{sealde} unchanged with \textbf{ðū}, misses the respective conjugation.
% mdparagraph:P197:end

\end{minipage}\par\addvspace{4pt}

\par\begin{minipage}{\linewidth}\raggedright
% mdparagraph:P198:begin
\textbf{X2.}
% mdparagraph:P198:end

% item: NEW-U05-X02
% key: KEY-NEW-U05-X02
% mdtable:T33:begin
\par\begingroup
\fontsize{10.5}{13.2}\selectfont
\setlength{\LTleft}{0pt}\setlength{\LTright}{\fill}
\noindent
\begin{tabular}{@{}>{\raggedright\arraybackslash}p{20.06mm}>{\raggedright\arraybackslash}p{20.06mm}>{\raggedright\arraybackslash}p{78.25mm}@{}}
\toprule
\textbf{Verb} & \textbf{Present 1sg} & \textbf{Past evidence in the supplied entry} \\
\midrule
swęrian & swęrie & swōr/\allowbreak{}swōron: strong vowel alternation, without a dental preterite suffix \\
lufian & lufie & lufode: weak dental preterite \\
\bottomrule
\end{tabular}
\par\endgroup

% mdtable:T33:end
% mdparagraph:P199:begin
An infinitive ending in \textbf{-ian} does not by itself identify a weak verb. \textbf{Swęrian} is strong VI despite its special present pattern; \textbf{lufian} is weak Class II. Their first-singular forms both end in \textbf{-ie}, so that one present ending alone also fails to distinguish the past systems. The infinitives and participles in the question are supplied citations; neither participial syntax nor unaided citation recall is required. [4]
% mdparagraph:P199:end

\end{minipage}\par\addvspace{4pt}

\par\begin{minipage}{\linewidth}\raggedright
% mdparagraph:P200:begin
\textbf{X3.}
% mdparagraph:P200:end

% item: NEW-U05-X03
% key: KEY-NEW-U05-X03
% mdtable:T34:begin
\par\begingroup
\fontsize{10.5}{13.2}\selectfont
\setlength{\LTleft}{0pt}\setlength{\LTright}{\fill}
\noindent
\begin{tabular}{@{}>{\raggedright\arraybackslash}p{20.58mm}>{\raggedright\arraybackslash}p{27.37mm}>{\raggedright\arraybackslash}p{43.94mm}>{\raggedright\arraybackslash}p{23.67mm}@{}}
\toprule
\textbf{Clause} & \textbf{Verb and strong class} & \textbf{Finite analysis} & \textbf{Meaning} \\
\midrule
Gē sēoð. & sēon, V & Present indicative, second plural & You, plural, see. \\
Hē sīeð. & sēoðan, II & Present indicative, third singular & He boils. \\
\bottomrule
\end{tabular}
\par\endgroup

% mdtable:T34:end
% mdparagraph:P201:begin
\textbf{Sēoð} is the contracted plural of “see.” \textbf{Sīeð}, with long \textbf{īe}, is the selected third-singular present of “boil.” The third-singular seeing form is instead \textbf{siehð}, with short \textbf{ie} and \textbf{h}. Do not identify a verb from an approximate visual resemblance or ignore \textbf{gē/\allowbreak{}hē}. These are form-identification examples within the indicative; no object or other complement needs to be invented. [3, 6]
% mdparagraph:P201:end

\end{minipage}\par\addvspace{4pt}

\clearpage
% mdparagraph:P202:begin
\section*{Source notes}
% mdparagraph:P202:end

\begingroup\fontsize{10}{12.8}\selectfont
% mdparagraph:P203:begin
[1] Cook, \textbf{A First Book in Old English}, second edition (1895), §§92–95, printed pp.54–55 /\allowbreak{} PDF70–71: four principal parts, finite indicative endings, plural-based second-singular preterites and the contracted \textbf{fōn} paradigm. The models with ordinary preceding subjects use the selected plural endings, without importing the separate postposed-pronoun \textbf{-e} condition.
% mdparagraph:P203:end

% mdparagraph:P204:begin
[2] Cook §105, printed p.63 /\allowbreak{} PDF79, with Wright and Wright, \textbf{Old English Grammar}, second edition (1914), §§503–504, pp.260–261 /\allowbreak{} PDF280–281. These supply Class IV and the particular \textbf{niman/\allowbreak{}cuman} alternatives.
% mdparagraph:P204:end

% mdparagraph:P205:begin
[3] Cook §106, printed pp.63–64 /\allowbreak{} PDF79–80, and Wright §§505–507, pp.261–262 /\allowbreak{} PDF281–282. Cook's consonant rules §§28–37, pp.22–25 /\allowbreak{} PDF38–41, support the stated dental combinations and contracted alternatives. Cook's dictionary \textbf{ge-sēon}, p.294 /\allowbreak{} PDF310, explicitly supports current display \textbf{gesēon}; removing the compound hyphen is declared normalization. Exact \textbf{geseah} remains short.
% mdparagraph:P205:end

% mdparagraph:P206:begin
[4] Cook §107, printed pp.64–66 /\allowbreak{} PDF80–82, and Wright §§508–510, pp.262–263 /\allowbreak{} PDF282–283. The present of \textbf{swęrian} is read with Cook §§116–117, pp.71–72 /\allowbreak{} PDF87–88. Complete finite cells apply the checked stems, endings and changes; they are not all individually printed lexical attestations.
% mdparagraph:P206:end

% mdparagraph:P207:begin
[5] Cook §§108–110, printed pp.66–67 /\allowbreak{} PDF82–83, and Wright §§511–519, pp.263–265 /\allowbreak{} PDF283–285. Wright §476, p.248 /\allowbreak{} PDF268, explicitly supports \textbf{fielþ} and the \textbf{x + ending → xt} result in \textbf{wiext}; \textbf{þ → ð} is the declared spelling convention in the classroom form \textbf{fielð}. Cook Exercises, vocabulary “grow,” p.50 /\allowbreak{} PDF58, and Cook's dictionary p.297 /\allowbreak{} PDF313 list \textbf{grōwan}. Its alternative past \textbf{grēowon} applies the \textbf{flōwan} pattern explicitly assigned to it in §109, p.67 /\allowbreak{} PDF83; it is not an individually printed past token on the dictionary page.
% mdparagraph:P207:end

% mdparagraph:P208:begin
[6] Cook §§101–104, printed pp.58–63 /\allowbreak{} PDF74–79, supply the Class I–III patterns used by the four new bridge verbs. Their specific meanings and classes are checked against Cook Exercises vocabulary and the lexical facsimiles. Bosworth–Toller, main dictionary (1898), \textbf{swimman}, p.957a, records the \textbf{o}-past alternatives; the lesson selects the Cook nasal-Class-III \textbf{a/\allowbreak{}u} series. The long \textbf{īe} of \textbf{flīehð/\allowbreak{}sīeð} must not be confused with the short \textbf{ie} of seeing \textbf{siehð}.
% mdparagraph:P208:end

% mdparagraph:P209:begin
[7] The bounded ordinary verb uses are individually checked in Bosworth–Toller facsimiles: \textbf{beran}, Main p.88a; \textbf{niman}, p.721a–b; \textbf{giefan}, p.474b; \textbf{etan}, p.259b and Supplement (1921) p.195a; \textbf{sittan}, p.879b; \textbf{licgan}, pp.636b–637a; \textbf{standan}, p.910b; \textbf{feallan}, p.271b; \textbf{healdan}, p.517b; \textbf{sēon}, p.864b; \textbf{ge-sēon}, p.439b; \textbf{flēogan}, p.291b; \textbf{scīnan}, p.833a; and \textbf{swimman}, p.957a. Food, birds and the other classroom combinations are declared new applications of these frames, not verbatim source quotations. Prior nominal and possessive paradigms remain those taught in Units 1–4.
% mdparagraph:P209:end

% mdparagraph:P210:begin
[8] Cook §156.k, printed p.92 /\allowbreak{} PDF108, explicitly illustrates \textbf{geswīcan} with the genitive of the work ceased from. The Sabbath passage, p.128 /\allowbreak{} PDF144, confirms the construction. Its English “from” does not require a separate Old English preposition in the selected frame.
% mdparagraph:P210:end

% mdparagraph:P211:begin
[9] Cook, \textbf{Exercises in Old English} (1895), Group I, Exercise 5, printed p.3 /\allowbreak{} PDF11, items 1–16. Original prompts are preserved. Item 11 has no note constraining the number of English \textbf{fish}; item 15's \textbf{swims} supplies that constraint for its own subject. Vocabulary “create,” p.44 /\allowbreak{} PDF52, and “see,” p.61 /\allowbreak{} PDF69, are read with the independently checked prefixed and unprefixed entries.
% mdparagraph:P211:end

% mdparagraph:P212:begin
[10] Cook, Reader I, \textbf{The Creation of the World}, first day: printed p.123 editorial body lines1–5 through p.124 line3 /\allowbreak{} PDF139–140; second day: p.124 lines4–10 /\allowbreak{} PDF140. On p.124, these line labels expand the printed marginal 5 and 10; p.123 counts are editorial. The 34 applicable numbered source notes are p.123 nn.1–18 and p.124 nn.1–16. Cook §53, p.37 /\allowbreak{} PDF53, supports weak feminine \textbf{heofone}; §87.d, p.52 /\allowbreak{} PDF68, the relative particle; and §173, p.97 /\allowbreak{} PDF113, the supplied naming construction with two accusatives. The first day is the unchanged earlier reading revisited; the second day is its complete new continuation.
% mdparagraph:P212:end

\endgroup
\end{document}
